\documentclass{studentnotes}

\setnotetitle{Theorem and Note-Box Compatibility Test}
\setnotecourse{LaTeX Lab}
\setnoteauthor{OT}
\setnotedate{4 August 2026}

\begin{document}

\makenotetitle

\section{Numbering and optional headings}

\begin{theorem}[Pythagorean theorem]
    \label{thm:pythagorean}
    For a right-angled triangle with legs \(a\) and \(b\), and
    hypotenuse \(c\),
    \[
        a^2+b^2=c^2.
    \]
\end{theorem}

\begin{theorem}
    \label{thm:unheaded}
    If \(n\) is even, then \(n^2\) is even.
\end{theorem}

\begin{definition}[Linear independence]
    \label{def:linear-independence}
    The vectors \(\mathbf{v}_1,\ldots,\mathbf{v}_n\) are linearly
    independent if
    \[
        c_1\mathbf{v}_1+\cdots+c_n\mathbf{v}_n=\mathbf{0}
    \]
    implies \(c_1=\cdots=c_n=0\).
\end{definition}

\begin{example}[Testing two vectors]
    \label{ex:independent-vectors}
    The vectors
    \[
        \begin{bmatrix}1\\0\end{bmatrix}
        \quad\text{and}\quad
        \begin{bmatrix}0\\1\end{bmatrix}
    \]
    are linearly independent.
\end{example}

The expected numbering in this section is:

\begin{itemize}
    \item the headed theorem: \(1.1\);
    \item the unheaded theorem: \(1.2\);
    \item the definition: \(1.1\);
    \item the example: \(1.1\).
\end{itemize}

The theorem, definition, and example counters are intentionally
independent under the current class definition.

\section{Section-based counter resets}

\begin{theorem}[Section-reset theorem]
    \label{thm:section-reset}
    For every real number \(x\),
    \[
        x^2\geq 0.
    \]
\end{theorem}

\begin{definition}[Absolute value]
    \label{def:absolute-value}
    The absolute value of a real number \(x\) is
    \[
        |x|=
        \begin{cases}
            x,  & x\geq 0,\\
            -x, & x<0.
        \end{cases}
    \]
\end{definition}

\begin{example}[Absolute value calculation]
    \label{ex:absolute-value}
    Since \(-4<0\), it follows that
    \[
        |-4|=4.
    \]
\end{example}

The theorem, definition, and example above should each be numbered
\(2.1\).

\section{Cross-reference behaviour}

Numeric references should produce the stored numbers:

\begin{itemize}
    \item Theorem~\ref{thm:pythagorean};
    \item Theorem~\ref{thm:unheaded};
    \item Definition~\ref{def:linear-independence};
    \item Example~\ref{ex:independent-vectors};
    \item Theorem~\ref{thm:section-reset};
    \item Definition~\ref{def:absolute-value};
    \item Example~\ref{ex:absolute-value}.
\end{itemize}

Automatic references should identify the corresponding environment:

\begin{itemize}
    \item \autoref{thm:pythagorean};
    \item \autoref{def:linear-independence};
    \item \autoref{ex:independent-vectors};
    \item \autoref{thm:section-reset}.
\end{itemize}

The optional headings should be available as descriptive references:

\begin{itemize}
    \item ``\nameref{thm:pythagorean}'';
    \item ``\nameref{def:linear-independence}'';
    \item ``\nameref{ex:independent-vectors}'';
    \item ``\nameref{thm:section-reset}''.
\end{itemize}

\section{Ordinary note-box rendering}

\begin{quicknote}
    This is an ordinary quick note. Its expected title is
    \textbf{Quick Note}, with the established yellow background and
    orange frame.
\end{quicknote}

\begin{personalnote}
    This is an ordinary personal note. Its expected title is
    \textbf{Personal Note}, with the established blue background and
    blue frame.
\end{personalnote}

\begin{importantnote}
    This is an ordinary important note. Its expected title is
    \textbf{Important}, with the established green background and
    green frame.
\end{importantnote}

The three boxes should retain their fixed titles, colours, sharp
corners, and established border thickness.

\clearpage
\section{Quick-note page-boundary test}

The following vertical space deliberately leaves too little room for
the complete note box at the bottom of this page.

\rule{0pt}{0.78\textheight}

\begin{quicknote}
    This quick note tests page-boundary behaviour. The complete box
    should move to the following page when insufficient space remains.
    It should not overlap the footer, become clipped, or lose its title
    and frame.
\end{quicknote}

\clearpage
\section{Personal-note page-boundary test}

The following vertical space deliberately leaves too little room for
the complete note box at the bottom of this page.

\rule{0pt}{0.78\textheight}

\begin{personalnote}
    This personal note tests page-boundary behaviour. The complete box
    should move to the following page when insufficient space remains.
    It should not overlap the footer, become clipped, or lose its title
    and frame.
\end{personalnote}

\clearpage
\section{Important-note page-boundary test}

The following vertical space deliberately leaves too little room for
the complete note box at the bottom of this page.

\rule{0pt}{0.78\textheight}

\begin{importantnote}
    This important note tests page-boundary behaviour. The complete box
    should move to the following page when insufficient space remains.
    It should not overlap the footer, become clipped, or lose its title
    and frame.
\end{importantnote}

\end{document}